%% ============================================================
%% SECTION: Introduction
%% ============================================================

\section{Introduction}

Brain tumor \cite{deangelis2001brain} encompasses various abnormal growths called neoplasms with distinct biology, prognosis, and treatment, often referred to as intracranial neoplasms. Modern medical science has advanced significantly in the understanding of the human brain, now using technologies like Magnetic Resonance Imaging (MRI) \cite{hashemi2012mri,brown2011mri} for non-invasive visualization and diagnosis of life-threatening diseases such as brain tumors. Magnetic Resonance Imaging (MRI) is a highly promising imaging modality that provides multicontrast information and fine anatomical detail for accurate tumor segmentation and precise boundary delineation, while also supporting tumor growth prediction, early prognosis, treatment optimization, and improved evaluation of patient outcomes \cite{amin2020distinctive}. Despite the advantages of MRI, brain tumor segmentation remains an inherently complex and challenging task due to the complex nature of certain tumor types, such as gliomas and glioblastomas. Gliomas and glioblastomas \cite{bahadure2017image} are difficult to accurately segment due to their diffuse boundaries, irregular shapes, and poor contrast with surrounding tissue, along with the wide variability in their appearance in different regions of the brain and imaging protocols. In addition to that, Glioblastomas exhibit infiltrative growth, which makes their margins indistinct from normal tissue. To improve delineation, multiple MRI modalities such as T1, T1 contrast enhanced (T1ce), T2, PD, dMRI, and fluid attenuated inversion recovery (FLAIR) are utilized, as complementary contrasts provide quasi-unique tissue signatures \cite{li2019multi, islam2021glioblastoma, rathore2018radiomic, hasan2016segmentation}.

The lack of standardization in the intensities of the MRI voxels further complicates the segmentation, as identical tumor tissues can present differently depending on the strength of the scanner and the acquisition settings. Consequently, automated annotation of volumetric MRI images has become an essential deep learning endeavor in supporting healthcare professionals to simplify and unify the process of labeling intricate 3D medical images for precise and consistent analysis. Conventional automatic brain tumor segmentation techniques rely predominantly on handcrafted feature engineering \cite{menze2014multimodal}, following a traditional machine learning framework where feature extraction and classifier training are performed independently. In contrast, deep neural networks enable end-to-end learning of hierarchical, task-specific feature representations directly from multimodal MRI data, facilitating more effective brain tumor segmentation \cite{bengio2013representation}.

{\sloppy Accurate and efficient brain tumor segmentation from multimodal MRI remains a challenge due to the substantial variability in tumor appearance, size, location, and morphology. Mohammad Havaei et al.~\cite{havaei2017brain} proposed a fully automatic deep convolutional neural network framework that integrates local and global contextual features, class-imbalance-aware training, and cascaded refinement to achieve competitive BraTS segmentation performance with reduced computational complexity and inference time. El-Sayed A et al.~\cite{el2014computer} proposed a hybrid machine learning framework that combines pulse-coupled neural network-based segmentation, wavelet feature extraction, principal component analysis, and neural network classification, achieving high accuracy in distinguishing normal and abnormal brain MRI scans while demonstrating robustness and computational efficiency. In addition, their comprehensive review of contemporary CAD-based brain MRI segmentation and classification techniques \cite{maitra2008hybrid, el2010hybrid, cheng2010automated, cherkassky2007learning} provided valuable insights into existing challenges and future research directions.\par}

Despite the promising results reported by existing CAD systems and deep learning-based approaches, accurate segmentation of heterogeneous brain tumor subregions from volumetric MRI data remains a challenging problem. The limitations of handcrafted feature-based methods and the growing demand for exploiting three-dimensional spatial context have encouraged the development of 3D MRI segmentation frameworks capable of learning more discriminative tumor representations. Motivated by these observations, this work proposes a novel 3D deep learning-based brain tumor segmentation framework to improve segmentation accuracy and robustness while maintaining computational efficiency.

\vspace*{-1\baselineskip}
